%% sections/model.tex
\section{Model}
\label{sec:model}

\subsection{Asset and fundamental value}

The asset pays an iid per-period dividend
$d_t\sim\text{Uniform}\{0,20\}$ with $\E[d_t]=10$ and lives for
$T=10$ periods per round.  The risk-neutral fundamental at the start
of period~$t$ is
\begin{equation}
\FV_t = 10\,(T-t+1),
\label{eq:fv}
\end{equation}
declining from $\FV_1=100$ to $\FV_T=10$.  Terminal value is zero.
The fundamental resets at every round boundary.  A session consists of
$R=4$ rounds; each period is discretised into $K=18$ ticks, giving
720 ticks per session.  The asset, dividend, and session structure
reproduce DLM exactly.

\subsection{Market institution}

Trade takes place in a continuous double auction with price-time
priority.  Bids are sorted descending by price, asks ascending; ties
break by arrival time.  Each tick, every agent submits at most one
order.  A crossing pair matches at the resting price; self-matches are
prevented.  At the end of each period, all outstanding orders are
withdrawn and the dividend is drawn and credited.

\subsection{Agent types}

Five agent types populate the market.  All observe the best bid~$b_t$,
best ask~$a_t$, last trade price~$p_t$, and fundamental~$\FV_t$.

\paragraph{Fundamentalist (F)}
Belief: $\beta_t = \FV_t(1+0.04(r{-}0.5))$, a $\pm2\%$ perturbation.
Buys at~$a_t$ when $a_t<0.98\beta_t$; sells at~$b_t$ when
$b_t>1.02\beta_t$; otherwise posts passive quotes near $\beta_t$ with
probability~$0.45$ or holds.

\paragraph{Trend follower (T)}
Estimates a slope from the last six in-round trade prices and
projects three ticks ahead.  Chases momentum when slope$\,{>}\,0$;
panic-sells when slope$\,{<}\,0$; liquidates excess inventory in
flat markets.  This is the positive-feedback channel
\citep{delong1990positive}.

\paragraph{Zero-intelligence (R)}
Posts $\text{bid}\sim\text{Uniform}[0.7\FV_t,1.1\FV_t]$ and
$\text{ask}\sim\text{Uniform}[0.9\FV_t,1.3\FV_t]$, following
\citet{gode1993allocative}.  Ensures the book is never empty.

\paragraph{DLM trader (D)}
Used only in strict-DLM mode ($N{=}6$).  A two-branch decision rule
is gated on the experience counter~$k_i$, which starts at zero and is
incremented at each round boundary:
\begin{equation}
\text{decide}() =
\begin{cases}
\text{inexperienced rule}, & k_i = 0,\\
\text{experienced rule}, & k_i \ge 1.
\end{cases}
\label{eq:dlm-gate}
\end{equation}

\emph{Inexperienced} ($k_i{=}0$).---The agent chases price momentum
with no fundamental anchor: when slope$\,{>}\,0$ it bids at
$p_t(1.02{+}0.04r)$; when slope$\,{<}\,0$ it asks at
$p_t(0.98{-}0.04r)$; otherwise it posts noisy quotes around~$p_t$.
This branch generates the bubble.

\emph{Experienced} ($k_i{\ge}1$).---The agent anchors to a
horizon-discounted belief
$\beta_t^{\text{exp}} = \FV_t\bigl(0.92 + 0.06\,(T{-}t{+}1)/T\bigr)$
and trades accordingly: it liquidates inventory when three or fewer
periods remain, refuses overpaying in late periods, crosses the book
only when the spread offers a discount or premium relative to
$\beta_t^{\text{exp}}$, and posts passive quotes near the belief with
probability~$0.35$.  This branch kills the bubble.

\paragraph{Utility agent (U)}
Follows \citet{lopezlira2025llm}.  On each tick the prior is
\begin{equation}
V_i^{\text{prior}} = \FV_t(1+b_i+\varepsilon_{it}),
\quad
b_i\in\{-0.15,\,0,\,0.15\},
\quad
\varepsilon_{it}\sim\text{Uniform}[-0.03,0.03],
\label{eq:util-prior}
\end{equation}
where $b_i$ is a persistent optimistic, unbiased, or pessimistic bias.
The agent scores each action $a\in\{\text{hold, cross bid, cross ask,
passive bid, passive ask}\}$ by expected utility:
\begin{equation}
\EU(a) = p_f\,U(w_1) + (1{-}p_f)\,U(w_0),
\qquad
a^\star = \arg\max_a\EU(a),
\label{eq:eu}
\end{equation}
with constant fill probability $p_f{=}0.30$ and risk-preference
functionals normalised so $U(w_0){=}1$:
\begin{equation}
\Uloving(w) = (w/w_0)^2,\qquad
\Uneutral(w) = w/w_0,\qquad
\Uaverse(w) = \sqrt{w/w_0}.
\label{eq:risk}
\end{equation}

\subsection{Communication and trust}

Each period, utility agents broadcast private valuations.  Three
deception modes govern the report: \emph{honest}
($\hat v_i = v_i(1{\pm}0.01)$), \emph{biased}
($\hat v_i = v_i(1{\pm}0.10)$, fixed sign), and \emph{deceptive}
($\hat v_i = 0.82v_i$ when under-invested, $1.18v_i$ when
over-invested).  Pairwise trust $\trust_{rs}\in[0,1]$ is initialised
at~$0.5$ and updated by exponential moving average against the
period's VWAP:
\begin{equation}
\trust_{rs} \leftarrow
(1{-}\alpha)\trust_{rs}
+ \alpha\max\!\bigl(0,\,1{-}|\hat v_s{-}\VWAP_t|/\VWAP_t\bigr),
\qquad \alpha=0.30.
\label{eq:trust}
\end{equation}
Trust accumulates across rounds; all other per-round state resets.

\subsection{Belief-update plans}

The period-boundary belief update maps the agent's prior and received
peer claims $\msgs_i$ into a posterior $V_i^{\text{post}}$ that
becomes next period's prior.  Three plans hold the upstream (prior,
messages, trust) and downstream (EU scoring, order matching, dividend
credit) stages fixed.

\paragraph{Plan I (algorithmic blend)}
The posterior is the experience-weighted convex combination of
equation~\eqref{eq:plan1}: a novice ($k_i{=}0$) weights peers at
$40\%$; a three-round veteran ($k_i{\ge}3$) at~$10\%$.  If no
messages arrive, the prior is unchanged.  Plan~I is deterministic
under the seeded PRNG.

\paragraph{Plan II (LLM direct action with utility expression)}
At each period boundary, one language-model call per utility agent
bypasses equation~\eqref{eq:plan1} entirely.  The structured prompt
encodes market rules, the agent's status, decision-making principles,
role-specific guidance (risk-type and experience-level), and the
\emph{closed-form} risk-utility expression:
$(w/w_0)^2$ for risk-loving, $w/w_0$ for neutral, $\sqrt{w/w_0}$ for
averse.  The LLM returns a discrete action from the set
$\{\text{BUY\_NOW}, \text{SELL\_NOW}, \text{BID\_1}, \text{BID\_3},
\text{ASK\_1}, \text{ASK\_3}, \text{HOLD}\}$,
which the engine translates into an order using the current book state.
A failed or invalid action degrades silently to Plan~I's EU evaluation.

\paragraph{Plan III (LLM direct action with risk label only)}
Identical to Plan~II except the functional form is deleted from the
prompt; only the category \emph{risk-loving}, \emph{risk-neutral}, or
\emph{risk-averse} is named.  Every other input and the seven-action
output set are identical.

If Plans~II and~III yield different actions, the difference is
attributable to prompt structure alone.  If a run's LLM calls all
fail, the trajectory is identical to the matched Plan~I run under the
same seed.
